\documentclass[aspectratio=169,14pt]{beamer}
\usepackage[utf8]{inputenc}
\usepackage[T1]{fontenc}
\usetheme{Madrid}
\setbeamertemplate{navigation symbols}{}
\title{Beamer Skills Library}
\subtitle{From research notes to presentation-ready slides}
\begin{document}
\begin{frame}[plain]
\titlepage
\end{frame}

\begin{frame}{Slide production follows a staged workflow}
\begin{itemize}
  \item Draft project progress as editable slides
  \item Convert structured content into Beamer source
  \item Apply academic presentation rules before delivery
\end{itemize}
\end{frame}

\begin{frame}{Three skills cover the authoring workflow}
\begin{itemize}
  \item beamer-format: content and layout standards
  \item beamer-live-draft: interactive HTML drafting and transpilation
  \item verify-tikz-layout: rendered evidence and visual inspection
\end{itemize}
\end{frame}

\begin{frame}{Progress drafts stay simple and convertible}
\begin{itemize}
  \item One HTML slide maps to one future Beamer frame
  \item Title slides and title-plus-bullets slides are supported
  \item JSON provides a stable handoff between editing and generation
\end{itemize}
\end{frame}

\begin{frame}{The generator preserves content and flags risks}
\begin{itemize}
  \item Escapes LaTeX special characters automatically
  \item Warns about empty titles and more than five bullets
  \item Supports clean, metropolis, and rochester templates
\end{itemize}
\end{frame}

\begin{frame}{Academic rules keep the deck audience-friendly}
\begin{itemize}
  \item One point per slide with conclusion-driven titles
  \item Readable fonts, balanced whitespace, and minimal tables
  \item Main results first; full details move to the appendix
\end{itemize}
\end{frame}

\begin{frame}{TikZ verification is a separate visual gate}
\begin{itemize}
  \item Compile and render the relevant pages or overlays
  \item Inspect every required image for overlap and clipping
  \item Repair layout without changing analytical meaning
\end{itemize}
\end{frame}

\begin{frame}{Notes become a reproducible PDF workflow}
\begin{itemize}
  \item Edit progress-draft.html and export progress.json
  \item Run generate.py to produce slides.tex and optionally slides.pdf
  \item Review the final deck with the publication and visual checks
\end{itemize}
\end{frame}
\end{document}
